\documentclass[11pt]{article}
\usepackage[T1]{fontenc}
\usepackage[utf8]{inputenc}
\usepackage{lmodern,microtype,amsmath,amssymb,amsthm,mathtools}
\usepackage[margin=1in]{geometry}
\usepackage{booktabs,longtable,array,enumitem}
\usepackage[colorlinks=true,linkcolor=blue,citecolor=blue,urlcolor=blue]{hyperref}
\usepackage{fancyhdr}
\setlength{\headheight}{14pt}
\pagestyle{fancy}\fancyhf{}
\fancyhead[L]{\small GCT\_GINEQ}
\fancyhead[R]{\small Norms, certificates, and tensor rank}
\fancyfoot[C]{\thepage}
\setlist{nosep}\setlength{\emergencystretch}{2em}
\newcommand{\R}{\mathbb R}\newcommand{\C}{\mathbb C}
\newcommand{\Z}{\mathbb Z}\newcommand{\Q}{\mathbb Q}
\newcommand{\F}{\mathbb F}
\newcommand{\BR}{\underline{\mathrm R}}\newcommand{\TR}{\mathrm R}
\newcommand{\SR}{\mathrm R_{\mathrm{sym}}}
\newcommand{\BSR}{\underline{\mathrm R}_{\mathrm{sym}}}
\newcommand{\HS}[1]{\lVert #1\rVert_{\mathrm F}}
\newcommand{\op}[1]{\lVert #1\rVert_{\mathrm{op}}}
\newcommand{\nuc}[1]{\lVert #1\rVert_*}
\newcommand{\eps}{\varepsilon}
\newcommand{\GL}{\mathrm{GL}}\newcommand{\Sym}{\operatorname{Sym}}
\newcommand{\Spec}{\operatorname{Spec}}\newcommand{\rank}{\operatorname{rank}}
\newcommand{\sgn}{\operatorname{sgn}}\newcommand{\diag}{\operatorname{Diag}}
\newcommand{\tr}{\operatorname{tr}}\newcommand{\id}{\operatorname{id}}
\newcommand{\Span}{\operatorname{span}}\newcommand{\Gr}{\operatorname{Gr}}
\newcommand{\Hilb}{\operatorname{Hilb}}\newcommand{\coker}{\operatorname{coker}}
\newcommand{\calO}{\mathcal O}\newcommand{\CP}{\mathbb P}
\newcommand{\abs}[1]{\lvert#1\rvert}
\newcommand{\ceil}[1]{\left\lceil#1\right\rceil}
\newenvironment{finding}[2]{\par\medskip\noindent\textbf{Finding #1 (#2).}\ \itshape}{\par\medskip}
\newenvironment{background}[2]{\par\medskip\noindent\textbf{Background #1 (#2).}\ }{\par\medskip}
\title{\textbf{Grothendieck's inequality and tensor norms\\against rank and border rank}\\[5pt]
\large Exact bounds, semidefinite certificates, and the algebraic-geometric dictionary}
\author{Research report for GCT\_GINEQ\\\small Run GCT\_GINEQ-2026-11-20-R1}
\date{20--21 November 2026\\\small Sources checked in this run; local date Europe/Berlin}
\begin{document}
\maketitle
\begin{abstract}
Projective and injective norms give exact information about matrix multiplication,
determinant, and permanent tensors, but matrix rank inequalities do not transfer
unchanged to tensor rank. We prove that none of three familiar full-tensor norm
ratios is bounded by any fixed constant times rank, even on square matrix
multiplication tensors. Valid support-dimension bounds are dominated by mode
flattenings. Matrix nuclear/operator certificates applied to a linear flattening,
including a Koszul flattening, do certify border rank; with arbitrary invertible
matrix preconditioning their optimal strength is exactly that of the same
flattening. Explicit singular spectra give $n^2$ for matrix multiplication,
$\binom n{\lfloor n/2\rfloor}$ for row determinant and permanent tensors, and
$\binom n{\lfloor n/2\rfloor}^2$ for their symmetric polynomial tensors.
A uniform injective-norm bound on the whole flattening image,
combined with border substitution, improves these bounds: it gives
$n^2+n-1$ for matrix multiplication and explicit additive
improvements for the determinant and permanent families. Coordinate
restrictions further reduce the residual-image problem exactly to
the border rank of a smaller member of the same family.
A dimension-sensitive Grothendieck gap gives a valid rank certificate, but the
natural flattenings here have gap one. We distinguish these deductions from
stronger known Koszul results, the Unique Games threshold, and recent numerical
preprints. A separate part gives the scheme, group-scheme, exterior-power,
Hilbert-scheme, and $K_0$ dictionary, and the precise limits of a translation
to weights or motives. No claim of a new rank record, exact Grothendieck
constant, or separation of complexity classes is made.
\end{abstract}
\clearpage
{\small\tableofcontents}
\newpage

\section*{Reading and verification conventions}
\addcontentsline{toc}{section}{Reading and verification conventions}
A \emph{finding} is a deduction formulated and proved in this run, not a
claim of historical priority. Known results are labelled background and
cited by statement or equation. The proofs are ordinary mathematical proofs,
checked by model rereading and, where recorded, separate reconstruction or
exact finite computations. They are not Lean proofs and have not been
externally refereed. \texttt{CHECK\_EVENTS.tsv} records the actual checks.
The model identifier is GPT-6 Astra Pro; a reasoning-effort setting was not
exposed and is not invented. No result of Alp\"oge or Masullo, and no
statement from an unseen earlier run, is used.

There are three different problems: row tensors have $n$ factors of
dimension $n$; the determinant/permanent polynomials give symmetric tensors
with $n$ factors of dimension $n^2$; the determinant--padded-permanent
orbit-closure problem is a further problem. These distinctions are essential.

\part{Checked analytic and algebraic findings}
\section{Norms, fields, and tensor conventions}
Let $V_1,\ldots,V_d$ be finite-dimensional real or complex Hilbert spaces
with specified standard metrics. Define
\[
 \eps(T)=\max_{\|x_i\|=1}
       \abs{\langle T,x_1\otimes\cdots\otimes x_d\rangle},\qquad
 \pi(T)=\inf_{T=\sum_\alpha\bigotimes_i x_{i\alpha}}
             \sum_\alpha\prod_i\|x_{i\alpha}\|.
\]
These are the Hilbert injective/spectral and projective/nuclear norms.
The complex pairing is Hermitian; absolute values make the equivalent
multilinear coefficient convention harmless. Nuclear norm is not a count
of terms, nor an average of the nuclear norms of matrix flattenings.

\begin{background}{B-01}{duality and symmetric norms}
The norms $\pi,\eps$ are dual, the nuclear infimum is attained, and
\begin{equation}\label{eq:duality}
 \abs{\langle Z,T\rangle}\leq\eps(Z)\pi(T),\qquad
 \HS T^2\leq\pi(T)\eps(T).
\end{equation}
For a symmetric tensor the injective norm is attained on a repeated unit
vector, and the nuclear norm admits a symmetric attaining decomposition.
Locators: Friedland--Lim \cite{FL}, \S2, equations (5)--(8) and the
following duality argument; Proposition 3.1; \S5, equation (23),
Theorem 5.1, Corollaries 5.2--5.3.
\end{background}

Write $\TR$ for CP rank. Over $\C$, $\BR$ is Zariski border rank,
equivalently Euclidean border rank. Over $\R$ we use Euclidean border
rank unless stated otherwise. Polynomial minor certificates also bound
complex border rank of a real tensor. Symmetric/Waring rank and border
rank, $\SR,\BSR$, allow scalar multiples of powers only. No rank version
of Comon's conjecture is assumed.
For $T\ne0$ put
\begin{equation}\label{eq:ratios}
 a(T)=\frac{\HS T^2}{\eps(T)^2},\qquad
 b(T)=\frac{\pi(T)^2}{\HS T^2},\qquad
 c(T)=\frac{\pi(T)}{\eps(T)}.
\end{equation}
Duality gives $a\leq c\leq b$. For matrices each is bounded by matrix
rank; this assertion is false for tensors.

With matrix units orthonormal, define
\[
 M_n=\sum_{i,j,k=1}^n E_{ij}\otimes E_{jk}\otimes E_{ki},\quad
 D_n=\sum_{\sigma\in S_n}\sgn(\sigma)\bigotimes_{i=1}^n e_{\sigma(i)},
 \quad
 P_n=\sum_{\sigma\in S_n}\bigotimes_{i=1}^n e_{\sigma(i)}.
\]
The last two are the \emph{row tensors}. Put $E=\F^{n\times n}$.
Let $\mathcal D_n,\mathcal P_n\in\Sym^nE$ have diagonal polynomials
$\det X,\operatorname{per}X$. Their ordered-basis coefficients are
$\pm1/n!$ and $1/n!$ on the $(n!)^2$ ordered perfect matchings of
matrix units, so their squared Hilbert norms are one. This tensor
normalization is in characteristic zero; Part II treats the integral form.

\begin{background}{B-02}{known exact norms}
Derksen \cite{Der}, Corollary 1.10, gives
$\pi(M_n)=n^3,\eps(M_n)=1$; Theorems 1.14--1.15 give
$\pi(P_n)=n^{n/2},\pi(D_n)=n!$. The former uses the sharp
permanent inequality of Carlen--Lieb--Loss \cite{CLL}, Theorem 1.1,
equation (1.1). The same values hold over both fields: real
decompositions provide the upper bounds and the complex spectral
inequalities give matching lower bounds. Coordinates: Derksen
\S1.5, equation (4), and \S1.7.
\end{background}

\begin{table}[ht]\centering\small
\begin{tabular}{lllll}\toprule
Tensor & order & $\HS T^2$ & $\eps(T)$ & $\pi(T)$\\\midrule
$M_n$ & $3$ & $n^3$ & $1$ & $n^3$\\
$D_n$ & $n$ & $n!$ & $1$ & $n!$\\
$P_n$ & $n$ & $n!$ & $n!/n^{n/2}$ & $n^{n/2}$\\
$\mathcal D_n$ & $n$ & $1$ & $n^{-n/2}$ & $n^{n/2}$\\
$\mathcal P_n$ & $n$ & $1$ & $n!/n^n$ & $n^n/n!$\\\bottomrule
\end{tabular}
\caption{Exact norms over both fields. The first three rows use B-02;
the last two are proved in F-04.}\label{tab:norms}
\end{table}

\section{Valid norm-to-rank conversions}
\begin{finding}{F-01}{support bounds and their limitation}
Let $d\geq2$, let $q_i$ be the mode-$i$ flattening rank of $T\ne0$,
and put $Q_T=\min_j\prod_{i\ne j}q_i$. Then
\begin{equation}\label{eq:support}
 \pi(T)\leq\sqrt{Q_T}\HS T,\qquad
 \eps(T)\geq\HS T/\sqrt{Q_T}.
\end{equation}
Consequently
\begin{equation}\label{eq:support-rank}
 \BR(T)\geq\ceil{b(T)^{1/(d-1)}}.
\end{equation}
The analogous bounds with $a,c$ are valid but weaker. Each is at most
$\max_iq_i$, the mode-flattening bound.
\end{finding}
\begin{proof}
The image of the mode-$i$ contraction is a support $E_i\subset V_i$
of dimension $q_i$. Projection to $E_i^\perp$ annihilates $T$, since
all contractions in the other factors vanish there. The commuting
orthogonal projections therefore show $T\in\bigotimes_i E_i$.
Projecting any decomposition, or test vectors, shows that both norms
computed in these smaller spaces equal the original norms.
Fix $j$ and orthonormal bases in the other supports. Expansion gives
$T=\sum_\alpha e_{\alpha_1}\otimes\cdots\otimes v_\alpha
\otimes\cdots\otimes e_{\alpha_d}$, with $P_j=\prod_{i\ne j}q_i$
summands. Hence
\[
 \pi(T)\leq\sum_\alpha\|v_\alpha\|
 \leq\sqrt{P_j}\Big(\sum_\alpha\|v_\alpha\|^2\Big)^{1/2}
 =\sqrt{P_j}\HS T.
\]
Testing against the basis vectors and $v_\alpha/\|v_\alpha\|$
gives $\eps(T)\geq\|v_\alpha\|$, and thus
$\HS T^2\leq P_j\eps(T)^2$. Minimize over $j$.
It follows that $a,b,c\leq Q_T$.
If $\BR(T)\leq r$, the $(r+1)$-minors of every mode flattening
vanish, since they vanish on every sum of $r$ pure tensors.
Thus $q_i\leq r$, so $Q_T\leq r^{d-1}$.
Finally $Q_T\leq(\max_iq_i)^{d-1}$; taking ceilings cannot
exceed the integer $\max_iq_i$.
\end{proof}
For a real tensor this particular bound with its real nuclear norm
also bounds complex border rank: real support expansion proves the
norm estimate, and real/complex mode ranks coincide. This is not a
general permission to substitute real norms in complex dual certificates.

\begin{finding}{F-02}{linear-flattening norm certificates and saturation}
Let $\Phi:\bigotimes_i V_i\to\operatorname{Mat}_{u,v}(\F)$ be linear
and let the integer $c_\Phi\geq1$ bound the rank of $\Phi$ on every
pure tensor. Put $A=\Phi(T)\ne0$. Certificates
$0\leq L\leq\nuc A$ and $U\geq\op A>0$ imply
\begin{equation}\label{eq:matrix-cert}
\BR(T)\geq\max\left\{
\ceil{\frac{L^2}{c_\Phi\HS A^2}},
\ceil{\frac{\HS A^2}{c_\Phi U^2}},
\ceil{\frac L{c_\Phi U}}\right\}.
\end{equation}
Each bound is at most $\ceil{\rank A/c_\Phi}$.
With exact norms and arbitrary invertible matrix preconditioning,
the best such bound equals $\ceil{\rank A/c_\Phi}$.
\end{finding}
\begin{proof}
Rank subadditivity and the vanishing of $(c_\Phi r+1)$-minors give
$\rank A\leq c_\Phi\BR(T)$. For its $s$ nonzero singular values $s_i$,
\[
 \frac{(\sum s_i)^2}{\sum s_i^2}\leq s,\qquad
 \frac{\sum s_i^2}{(\max s_i)^2}\leq s,\qquad
 \frac{\sum s_i}{\max s_i}\leq s.
\]
These prove the certificates and their ceiling. In a singular-vector
basis an invertible left preconditioner multiplying direction $i$
by $1/s_i$ and acting as the identity on the complementary subspace
makes every nonzero singular value one. Every ratio then equals $s$.
Preconditioning preserves ranks on all pure tensors, so $c_\Phi$
is unchanged.
\end{proof}
This identifies the strength of a specified linear map. It does not
assert that all uses of norms reduce to linear flattenings.

Finite semidefinite witnesses for F-02 are
\[
\begin{pmatrix}I_u&Z\\Z^*&I_v\end{pmatrix}\succeq0,\qquad
L=\operatorname{Re}\tr(Z^*A),\qquad
\begin{pmatrix}UI_u&A\\A^*&UI_v\end{pmatrix}\succeq0 .
\]
Schur complements give $\op Z\leq1$ and $\op A\leq U$; matrix duality
gives $L\leq\nuc A$. Rational Gram or $LDL^*$ factorizations,
or exact algebraic certificates with directed enclosures, suffice.
An approximate objective alone is not a certificate. These Hilbert
matrix SDPs require no Grothendieck constant.

\section{Explicit bounds for the tensor families}
\begin{finding}{F-03}{flat singular spectra}
A mode flattening of $M_n$ has $n^2$ singular values $\sqrt n$.
A $k\mid(n-k)$ flattening of $D_n$ or $P_n$ has $\binom nk$
singular values $\sqrt{k!(n-k)!}$.
A $k\mid(n-k)$ flattening of $\mathcal D_n$ or $\mathcal P_n$
has $\binom nk^2$ singular values $1/\binom nk$.
All other singular values vanish. Consequently
\begin{align}
 \BR(M_n)&\geq n^2,\label{eq:MMbound}\\
 \BR(D_n),\BR(P_n)&\geq\binom n{\lfloor n/2\rfloor},\label{eq:rowbound}\\
 \BR(\mathcal D_n),\BR(\mathcal P_n)&\geq
          \binom n{\lfloor n/2\rfloor}^{\!2}.\label{eq:polybounds}
\end{align}
The last bound also holds for symmetric border rank. The assertions
hold over both fields with the stated conventions.
\end{finding}
\begin{proof}
Every row of the first flattening of $M_n$ has $n$ unit entries.
Distinct rows have disjoint supports, because the remaining two
matrix-unit indices determine the first. Its row Gram matrix is $nI$.
For a row tensor, group a nonzero row by its underlying $k$-element
subset $I$, and its columns by $I^c$. The block is $k!$ by
$(n-k)!$. For $P_n$ its entries are one. For $D_n$ the sign is
a fixed shuffle sign times the sign of the row order and the sign
of the column order. Thus each block is an outer product of sign
vectors, with singular value $\sqrt{k!(n-k)!}$.
The blocks have disjoint row and column supports.

For a symmetric polynomial tensor a row consists of $k$ matrix units
with distinct row set $I$ and column set $J$. For each pair $(I,J)$
there are $(k!)^2$ such ordered lists, and $((n-k)!)^2$
complementary column lists. The permanent block is constant $1/n!$.
The determinant sign factors into signs of the partial matching,
the complementary matching, and a fixed shuffle sign depending on
$(I,J)$. Thus every block has rank one and singular value
$k!(n-k)!/n!=1/\binom nk$. There are $\binom nk^2$ blocks.
All three matrix norm ratios equal rank for a flat nonzero singular
spectrum. Apply F-02 with $c_\Phi=1$ and maximize at
$k=\lfloor n/2\rfloor$. A power is pure, so these minors also
vanish on the symmetric border-rank locus.
\end{proof}
These are Hilbert-normalized catalecticants. Differential conventions
can change basis lengths by factorials, so their singular values must
not be quoted without a normalization.

\begin{finding}{F-04}{symmetric polynomial norms and nuclear decompositions}
Over both fields,
\[
\eps(\mathcal D_n)=n^{-n/2},\quad\pi(\mathcal D_n)=n^{n/2},
\qquad
\eps(\mathcal P_n)=n!/n^n,\quad\pi(\mathcal P_n)=n^n/n!.
\]
Real symmetric nuclear decompositions exist with at most
$n!2^{n-1}$ and $4^{n-1}$ terms, respectively.
\end{finding}
\begin{proof}
By B-01 maximize the diagonal polynomials on $\HS X=1$.
The product of singular values gives
$|\det X|\leq n^{-n/2}$, attained at $I/\sqrt n$.
The CLL inequality followed by arithmetic-geometric mean for the
row norms gives
$|\operatorname{per}X|\leq(n!/n^{n/2})\prod_i\|\operatorname{row}_iX\|
\leq n!/n^n$, attained by the matrix all of whose entries are $1/n$.

Let $P_\sigma$ have its unit entry in row $i$ at column $\sigma(i)$.
For uniform $\sigma$ and independent signs $\eta_i$, let
$X_{\sigma,\eta}=\diag(\eta)P_\sigma/\sqrt n$, a real unit matrix.
Coefficientwise sign averaging gives
\begin{equation}\label{eq:det-average}
\mathbb E_{\sigma,\eta}
\left[\sgn(\sigma)\Big(\prod_i\eta_i\Big)X_{\sigma,\eta}^{\otimes n}\right]
=n^{-n/2}\mathcal D_n.
\end{equation}
Every row occupancy must be odd, hence exactly one since there are
$n$ factors. A surviving matching has a unique compatible $\sigma$
and coefficient $\sgn(\sigma)/(n!n^{n/2})$, proving the identity.
The total absolute weight is one, so $\pi(\mathcal D_n)\leq n^{n/2}$.
Duality and squared Hilbert norm one give equality.

For independent sign vectors $\eta,\theta$, put
$Y_{\eta,\theta}=(\eta_i\theta_j)_{ij}/n$. It has Hilbert norm one.
Sign averaging now forces every row and column occupancy to be one,
and gives
\begin{equation}\label{eq:perm-average}
\mathbb E_{\eta,\theta}
\left[\Big(\prod_i\eta_i\Big)\Big(\prod_j\theta_j\Big)
Y_{\eta,\theta}^{\otimes n}\right]
=\frac{n!}{n^n}\mathcal P_n.
\end{equation}
The surviving coefficient is $n^{-n}$. This proves the nuclear
upper bound; duality gives equality. Global negation of $\eta$
in \eqref{eq:det-average} leaves the weighted tensor unchanged.
In \eqref{eq:perm-average}, either independent global negation has
that property. The numbers of terms reduce to $n!2^{n-1}$ and
$2^{2n-2}$. These are decompositions by powers attaining the norm.
\end{proof}

\begin{center}\small
\begin{tabular}{lll}\toprule
Tensor & F-01 bound & F-03 bound\\\midrule
$M_n$ & $\ceil{n^{3/2}}$ & $n^2$\\
$D_n$ & $\ceil{(n!)^{1/(n-1)}}$ & $\binom n{\lfloor n/2\rfloor}$\\
$P_n$ & $\ceil{(n^n/n!)^{1/(n-1)}}$ & $\binom n{\lfloor n/2\rfloor}$\\
$\mathcal D_n$ & $\ceil{n^{n/(n-1)}}$ & $\binom n{\lfloor n/2\rfloor}^2$\\
$\mathcal P_n$ & $\ceil{(n^{2n}/(n!)^2)^{1/(n-1)}}$
 & $\binom n{\lfloor n/2\rfloor}^2$\\\bottomrule
\end{tabular}
\end{center}
Here $n\geq2$. By F-01 every full-norm conversion in the middle
column is dominated even by a mode flattening. Knowing the full
norm exactly has not improved rank through this conversion.

\section{Why matrix intuition fails}
\begin{finding}{F-05}{no fixed Grothendieck-factor repair}
No finite constant $C$ makes any one of $a(T),b(T),c(T)\leq C\TR(T)$
valid on all square matrix multiplication tensors. The border-rank
versions fail as well. The same failure occurs along each of the four
determinant/permanent families above; for the symmetric polynomial
families even Waring rank may replace rank on the right.
\end{finding}
\begin{proof}
All three ratios for $M_n$ equal $n^3$. Strassen's seven-term formula
gives $\TR(M_2)\leq7$ over both fields (Derksen \cite{Der},
\S1.5, equation (5)). Tensoring decompositions factorwise gives
$\TR(S\boxtimes T)\leq\TR(S)\TR(T)$.
The identification
$E_{ij}\otimes E_{ab}\mapsto E_{(i,a),(j,b)}$ in each factor proves
$M_n\boxtimes M_m=M_{nm}$. Consequently
\[
 \TR(M_{2^k})\leq7^k,\qquad
 a(M_{2^k})=b(M_{2^k})=c(M_{2^k})=8^k.
\]
The quotient $(8/7)^k$ is unbounded. This argument does not assume
multiplicativity of an arbitrary tensor norm.

For the other families, the real decompositions recorded in
Derksen \S1.7, equation (8), and \S8, p.~23 (the five-term
display and final unnumbered display before Remark 8.1), give
\[
 \TR(P_n)\leq2^{n-1},\qquad
 \TR(D_n)\leq(5/6)^{\lfloor n/3\rfloor}n!.
\]
Their three ratios are $n^n/n!$ and $n!$ respectively.
F-04 gives symmetric-rank upper bounds $n!2^{n-1}$ and $4^{n-1}$
for $\mathcal D_n,\mathcal P_n$, whose ratios are
$n^n$ and $(n^n/n!)^2$. Put $A_n=n^n/(2^{n-1}n!)$.
The binomial expansion gives, for $n\geq2$,
\[
 \frac{A_{n+1}}{A_n}
 =\frac{(1+1/n)^n}{2}
 \geq\frac{2+(n-1)/(2n)}2\geq\frac98.
\]
The quotients of norm ratios by the displayed rank upper bounds
are therefore at least $A_n$, $(6/5)^{\lfloor n/3\rfloor}$,
$A_n$, and $A_n^2$. All are unbounded.
Border rank is no larger than rank, so the border versions also fail.
\end{proof}
No fixed power of either Grothendieck constant repairs these
inequalities. Already $M_2$ refutes constant one; its powers refute
every fixed constant while keeping tensor order three.

\begin{finding}{F-06}{a degeneration and an exact rational counterexample}
Let $W=e_1\otimes e_1\otimes e_2+
e_1\otimes e_2\otimes e_1+
e_2\otimes e_1\otimes e_1$. Over both fields,
$\TR(W)=3$ and $\BR(W)=2$. The rational tensor
\[
 T_t=\frac{(e_1+t e_2)^{\otimes3}-e_1^{\otimes3}}t,\qquad t=1/100,
\]
has rank two and $a(T_t)>2$, hence $b(T_t),c(T_t)>2$.
\end{finding}
\begin{proof}
The slice pencil of $W$ is
$\left(\begin{smallmatrix}y&x\\x&0\end{smallmatrix}\right)$,
with determinant $-x^2$. Its rank-one members occupy only one
line. If $W$ had rank two, its mode rank two would force independent
first factors and would express this pencil as the span of two
independent rank-one matrices, a contradiction. Its three displayed
terms give rank three. For $t\ne0$, the two first factors and the
two squared remaining factors of $T_t$ are independent, so its rank
is exactly two. It tends to $W$; the mode rank of $W$ is two.
This proves the border assertion, including the real Euclidean case.

Friedland--Lim \cite{FL}, Lemma 6.2, equations (34)--(37), gives
after rescaling $\eps(W)=2/\sqrt3$ over both fields. Directly,
\[
 \HS{T_t}^2=3+3t^2+t^4\geq3,\quad
 \HS{T_t-W}=\sqrt{3t^2+t^4}\leq2t=1/50.
\]
Since $\eps\leq\HS{\cdot}$, the triangle inequality gives
$\eps(T_t)<6/5+1/50=61/50$. Hence
$a(T_t)>7500/3721>2$. Duality gives the other two inequalities.
\end{proof}
The same lemma gives $\pi_\R(W)=3$ and
$\pi_\C(W)=3\sqrt3/2$. Over $\C$ all three ratios at $W$ are $9/4$;
over $\R$ they are $9/4$, $3$, and $3\sqrt3/2$ in the order
$a,b,c$. The field distinction is material.

\begin{finding}{F-07}{bounded summands and unavoidable divergence}
Suppose $\eps(Z)\leq U$, $U>0$, $B>0$, $\delta\geq0$, and
$T'=\sum_{j=1}^r\lambda_j A_j$, with unit pure $A_j$,
$|\lambda_j|\leq B$, and $\HS{T-T'}\leq\delta$. Then
\begin{equation}\label{eq:bounded}
 r\geq
 \frac{(\abs{\langle Z,T\rangle}-\delta\HS Z)_+}{UB}.
\end{equation}
If $\TR(T)>r$, every sequence of decompositions of length at most
$r$ converging to $T$ has total term norm tending to infinity.
\end{finding}
\begin{proof}
We have $|\langle Z,T'\rangle|\leq U\sum_j|\lambda_j|\leq UrB$,
and Cauchy--Schwarz bounds the error pairing by $\delta\HS Z$.
For the second assertion, any bounded subsequence of total weights
bounds all the terms. The pure tensors of Hilbert norm at most
a fixed number form a compact set: they are a continuous image of
a bounded scalar interval times unit spheres. Pad by zero terms
and pass to a subsequence for each of the finitely many terms.
Their limits give a length-at-most-$r$ decomposition of $T$,
contradicting its rank.
\end{proof}
Thus a summand bound in \eqref{eq:bounded} certifies the specified
bounded-summand problem; it cannot be imposed in unrestricted border
rank. Also the convex hull of the full rank-one cone is the whole
tensor space: any finite sum becomes an average by scaling its terms.
Nuclear balls relax normalized atoms, not that unrestricted cone.

\section{A necessary source correction}
\begin{finding}{F-08}{correct flattening norm comparisons}
If a three-tensor $A$ has mode ranks $(q_1,q_2,q_3)$, put
$s=\min(q_2,q_3)$. Its first flattening satisfies
\[
 \nuc{A_{(1)}}\leq\pi(A)\leq\sqrt s\,\nuc{A_{(1)}},
 \qquad
 \eps(A)\leq\op{A_{(1)}}\leq\sqrt s\,\eps(A).
\]
The supplied Friedland--Lim v3 displays at equation (57) and immediately
before (60) cannot be used literally as their proof.
\end{finding}
\begin{proof}
Restrict to minimal supports as in F-01. Flattening a pure tensor
preserves its product norm, proving the first nuclear inequality.
Enlarging the set of spectral test vectors proves the first spectral
inequality. In an SVD
$A_{(1)}=\sum_\ell s_\ell x_\ell\otimes V_\ell$, every reshaped
$V_\ell$ has Hilbert norm one in a $q_2$ by $q_3$ matrix space,
and hence nuclear norm at most $\sqrt s$. Expanding these matrices
by SVD proves the nuclear upper bound. The contraction $A(x)$
for unit $x$ has operator norm at most $\eps(A)$, so for a
Hilbert-unit $V$,
$|\langle A,x\otimes V\rangle|\leq\eps(A)\nuc V
\leq\sqrt s\,\eps(A)$. This proves the other bound.

For the first source defect set $A=e\otimes I_m$, $m>1$, with $e$
unit. Its spectral norm is one, but the right side of (57), with
denominator $\|x\|\op M$, is at least $m$ at $x=e,M=I_m$.
The denominator must be $\|x\|\nuc M$.
The other displayed assertion,
$\nuc M\leq\sqrt{\rank M}\op M$, fails for $I_m$.
Its correct Frobenius version is
$\nuc M\leq\sqrt{\rank M}\HS M$.
\end{proof}
The standard conclusion is the form of \cite{FL}, Theorem 9.4,
needed here. These defect statements concern the supplied source and
its arXiv v3 PDF; no claim about every later typesetting is made.

\section{Grothendieck's inequality as a certificate}
\subsection{The Banach spaces and the directions of the bounds}
For a real $m$ by $n$ matrix define
\[
 Q(A)=\max_{x_i,y_j\in\{-1,1\}}\sum_{ij}a_{ij}x_i y_j,\qquad
 S(A)=\max_{\|u_i\|=\|v_j\|=1}\sum_{ij}a_{ij}\langle u_i,v_j\rangle .
\]
Absolute values give the same optima by negating one side. A Gram
matrix realizes the second maximum in dimension at most $m+n$.
Allowing norms at most one gives the same value, by padding with
private orthogonal coordinates.

\begin{background}{B-03}{the discrete Grothendieck inequality}
$Q(A)\leq S(A)\leq K_G^\R Q(A)$.
In tensor-norm language $Q$ is the injective norm on
$\ell_1^m\otimes\ell_1^n$, equivalently the bilinear-form norm on
$\ell_\infty^m\times\ell_\infty^n$, and $S=\gamma_2^*$.
The dual statement compares the projective norm on
$\ell_\infty^m\otimes\ell_\infty^n$ to $\gamma_2$.
These are not the Hilbert tensor norms used above.
Locators: Pisier \cite{Pis}, Theorem 2.4, (2.4)--(2.5);
Theorem 3.1, (3.13)--(3.14).
\end{background}

\begin{finding}{F-09}{directional SDP certificate}
Set $C_A=\frac12\left(\begin{smallmatrix}0&A\\A^T&0\end{smallmatrix}\right)$.
If $X\succeq0$ has diagonal one, and
$\diag(z)-C_A\succeq0$, then, with
$L=\langle C_A,X\rangle$ and $U=\sum_i z_i$,
\[
 L\leq S(A)\leq U,\qquad L/K\leq Q(A)\leq U
\]
for any proved $K\geq K_G^\R$.
\end{finding}
\begin{proof}
Factoring $X$ as a Gram matrix gives feasible unit vectors and value
$L$. For any feasible Gram matrix $Y$,
\[
 \langle C_A,Y\rangle
 =\langle\diag(z),Y\rangle-\langle\diag(z)-C_A,Y\rangle
 \leq\sum_i z_i,
\]
because the trace pairing of positive semidefinite matrices is
nonnegative. Apply B-03. No strong-duality assumption is required.
\end{proof}
In particular, division by $K$ supplies a \emph{lower} bound on
the sign norm. It does not turn a primal SDP value into a spectral
upper bound.

\begin{finding}{F-10}{the direct cube/Hilbert route and the natural targets}
For a real tensor $W$ and a coordinate flattening $F$,
\[
 \eps_\C(W)\leq\op{F(W)}\leq Q(F(W))\leq S(F(W)),
 \quad Q(F(W))\geq\HS W,\quad S(F(W))\geq\nuc{F(W)}.
\]
Using $S(F(T))$ directly as a spectral upper bound in F-01 gives
a rank lower bound at most one. A nuclear dual witness certified
solely by $S(F(W))\leq1$ gives at most $\op{F(T)}\leq\HS T$,
so this route in F-01 is also uninformative.
Moreover the natural flattenings in F-03 all satisfy
$Q(A)=S(A)=\sum_{ij}|a_{ij}|$.
\end{finding}
\begin{proof}
Grouping unit tensor factors gives unit matrix test vectors.
A real matrix has the same real and complex Hilbert operator norm;
real unit vectors lie in the coordinate cube, proving the first chain.
For independent random signs $x,y$,
$\mathbb E(x^TAy)^2=\sum a_{ij}^2$, so $Q(A)\geq\HS A$.
In a thin SVD $A=U\Sigma V^T$, the rows of $U,V$ have norm at most
one. Use these rows as SDP vectors and pad if necessary.
The objective is $\tr\Sigma$, proving $S(A)\geq\nuc A$.
Thus $S(F(W))\leq1$ implies $\nuc{F(W)}\leq1$, whence
\[
 |\langle T,W\rangle|
 =|\langle F(T),F(W)\rangle|\leq\op{F(T)}.
\]
This argument concerns real witnesses; for a real target such
witnesses also give valid complex dual bounds because the first
chain bounds $\eps_\C(W)$. The two F-01 conclusions follow.

For the last assertion, an entrywise nonnegative matrix has
$Q=S=\sum a_{ij}$: every vector inner product is at most one,
and all equal unit vectors attain equality. Row and column sign
switches preserve both optima. The $M_n$ and permanent matrices are
already nonnegative. In the determinant matrices, F-03 decomposed
the support into disjoint rank-one blocks with factored signs.
Each row and column belongs to one block, so all these signs can
be removed simultaneously. This proves the assertion.
\end{proof}
This is a limitation of a specified unscaled relaxation and of the
natural maps. It does not exclude preconditioning, tensor-specific
constructions, noncommutative inequalities, or higher SDP/SOS methods.

\begin{finding}{F-11}{a genuinely Grothendieck-based rank test}
Let $K_G^\R(s)$ mean the vector-dimension-$s$ Grothendieck constant
(Pisier \cite{Pis}, (4.4)). A nonzero real matrix of rank $q$ has an
optimal vector solution for $S(A)$ in dimension at most $q$.
Thus $S(A)\leq K_G^\R(q)Q(A)$.
If $L\leq S(A)$, $Q(A)\leq U$, and
$\kappa_s\geq K_G^\R(s)$ are independently certified, then
\begin{equation}\label{eq:rank-sensitive}
 L>\kappa_s U\quad\Longrightarrow\quad\rank A>s.
\end{equation}
For real $A=\Phi(T)$ with atom rank cost $c_\Phi$, the test with
$s=c_\Phi r$ certifies $\BR_\C(T)>r$ as well.
\end{finding}
\begin{proof}
Start with a maximizing vector family. Set $w_i=\sum_j a_{ij}v_j$.
Their span $E$ has dimension at most $q$, since it is the image
of the row space of $A$. Replacing $u_i$ by $w_i/\|w_i\|$ when
nonzero maximizes the objective with $v$ fixed and hence preserves
the global maximum. If $w_i=0$, choose any unit vector in $E$.
The space $E$ is nonzero because $A\ne0$ has positive optimum.
Now set $z_j=\sum_i a_{ij}u_i\in E$ and replace $v_j$ by its unit
direction, choosing an arbitrary unit vector of $E$ at a zero.
Again the global maximum is preserved. All vectors now lie in $E$.
The definition and monotonicity of $K_G^\R(s)$ prove
\eqref{eq:rank-sensitive}. Real and complex ranks of a real matrix
coincide; apply the minor argument of F-02.
\end{proof}
The universal $K_G$ cannot produce such a strict violation.
Nor can the same basic SDP upper bound serve as $U$:
then $L\leq S(A)\leq U$ and $\kappa_s\geq1$.
A separate bound on the sign optimum is necessary.
F-10 shows that this test gives no rank information on the
natural target matrices; even elsewhere it remains bounded by
the rank of the chosen flattening.

\subsection{A small exact certificate}
Let $A=\left(\begin{smallmatrix}1&1\\1&-1\end{smallmatrix}\right)$.
Use $u_1=(1,0),u_2=(0,1)$ and
$v_1=(3/5,4/5),v_2=(3/5,-4/5)$. All are rational unit vectors,
and their SDP objective is $14/5$.
For signs,
$Q(A)=|y_1+y_2|+|y_1-y_2|=2$.
Since $K_G^\R(1)=1$, F-11 proves $\rank A>1$.
This upper bound on $Q$ also has a finite SOS certificate.
Put $f=x_1(y_1+y_2)+x_2(y_1-y_2)$ and
\begin{align*}
 H_0={}&(x_1^2-1)(y_1+y_2)^2+(x_2^2-1)(y_1-y_2)^2\\
 &+(2+2x_1x_2)(y_1^2-1)+(2-2x_1x_2)(y_2^2-1).
\end{align*}
Expansion gives $f^2-4=H_0$ and
$2-f=\frac14(2-f)^2-\frac14H_0$.
Modulo the Boolean equations, the right side is a square.
The included exact-arithmetic script checks the Gram data and the
polynomial identity. This demonstrates a certificate, not a
uniform polynomial-size SOS theorem.

\section{Koszul flattenings and the precise method ceiling}
\begin{finding}{F-12}{Koszul cost, certificates, and dimension bound}
For $T\in A\otimes B\otimes C$, with dimensions $a,b,c$, let
\[
 K_p(T):B^*\otimes\Lambda^p A\longrightarrow
          \Lambda^{p+1}A\otimes C,\qquad 0\leq p<a,
\]
send $\beta\otimes\omega$ to
$\sum_\nu\beta(b_\nu)(a_\nu\wedge\omega)\otimes c_\nu$
when $T=\sum_\nu a_\nu\otimes b_\nu\otimes c_\nu$.
Its rank on a nonzero pure tensor is
$c_p=\binom{a-1}p$. Thus F-02 applies with $c_\Phi=c_p$, and
\[
 \frac{\rank K_p(T)}{c_p}
 \leq \min\left\{\frac{ba}{a-p},\frac{ca}{p+1}\right\}.
\]
Nuclear/operator SDP certificates on this map cannot exceed its
ordinary Koszul rank certificate.
\end{finding}
\begin{proof}
For $T=a_0\otimes b_0\otimes c_0\ne0$, contraction by $b_0$
has one-dimensional image, and tensoring by $c_0$ does not change rank.
Choose a basis with $a_0=e_1$. Wedge multiplication by $e_1$
annihilates basis wedges containing $e_1$ and maps the other
$\binom{a-1}p$ basis wedges to independent wedges.
Subadditivity and minors extend the bound to border rank.
The source and target dimensions are
$b\binom ap$ and $c\binom a{p+1}$; divide their minimum by
$\binom{a-1}p$. The last assertion is F-02.
\end{proof}
The construction works after a restriction in a factor as well.
It is a map of exterior-power bundles in Part II, not merely a
coordinate matrix.

\section{A stronger route: norms on an entire flattening image}
The preceding limitations do not exhaust norm-assisted arguments.
Border substitution introduces information about every tensor in a
flattening image, rather than only the rank of that flattening.
This leads to strictly stronger bounds than the basic bounds of F-03.

\begin{finding}{F-13}{image substitution with an injective-norm certificate}
Flatten $T$ into a map with image
$U\subseteq H_1\otimes\cdots\otimes H_b$, where $b\geq2$, and
write $q=\dim U>0$. Over $\C$,
\begin{equation}\label{eq:image-sub}
 \BR(T)\geq q+\min_{0\ne u\in U}\BR(u)-1.
\end{equation}
The analogous assertion holds for real Euclidean border rank.
Let $P_U$ be the Hilbert orthogonal projector and define
\begin{equation}\label{eq:alpha}
 \alpha(U)=\max_{\substack{z=x_1\otimes\cdots\otimes x_b\\\|x_i\|=1}}
                    \|P_Uz\|^2
           =\max_{\substack{u\in U\\\HS u=1}}\eps(u)^2.
\end{equation}
If a certified $0<\alpha_0\leq1$ bounds $\alpha(U)$ above, then
\begin{equation}\label{eq:image-norm}
 \BR(T)\geq q+\ceil{\alpha_0^{-1/(b-1)}}-1.
\end{equation}
Moreover $\alpha(U)$ is the injective norm of $P_U$ regarded as
a $2b$-tensor in the factors $H_i$ and their conjugate spaces.
\end{finding}
\begin{proof}
For \eqref{eq:image-sub}, let $\BR(T)=r$ and take rank-at-most-$r$
approximants $T_\nu\to T$. Choose $q$ fixed domain vectors whose
images under $T$ are independent. Their images under $T_\nu$
remain independent for large $\nu$, span $q$-planes $U_\nu$, and
$U_\nu\to U$ in the Grassmannian.
In a length-at-most-$r$ decomposition of $T_\nu$, let $S_\nu$
be the span of all the pure tensors in the remaining $b$ factors.
Then $U_\nu\subseteq\operatorname{im}T_\nu\subseteq S_\nu$ and
$q\leq s_\nu:=\dim S_\nu\leq r$.
Choose a basis of $S_\nu$ from those pure atoms, and take the span
$L_\nu$ of any $s_\nu-q+1$ of them. Dimension counting gives
$U_\nu\cap L_\nu\ne0$. A nonzero $u_\nu$ in that intersection
has rank at most $s_\nu-q+1\leq r-q+1$.
Normalize its Hilbert norm to one and pass to a convergent subsequence.
The limit $u$ has norm one, belongs to $U$, and has border rank
at most $r-q+1$. This proves \eqref{eq:image-sub}.
The same compactness proof works over $\R$ with Euclidean closure.
The atom span $S_\nu$, not the possibly smaller image of $T_\nu$,
is essential in this argument.

For \eqref{eq:alpha}, write
$\|P_Uz\|=\max_{\HS u=1,u\in U}|\langle u,z\rangle|$
and interchange compact maxima. If $\alpha(U)\leq\alpha_0$,
then $a(u)\geq1/\alpha_0$ for every $u\ne0$ in $U$.
F-01 on the remaining $b$ factors gives
$\BR(u)\geq\ceil{\alpha_0^{-1/(b-1)}}$.
For the projector tensor, its spectral test values are
$|\langle y,P_Ux\rangle|$ for unit pure $x,y$.
Positive semidefinite Cauchy--Schwarz bounds this by
$\sqrt{\langle x,P_Ux\rangle\langle y,P_Uy\rangle}\leq\alpha(U)$.
Taking $x=y$ at a maximizer gives equality.
\end{proof}
For three factors the border-substitution principle is known:
Landsberg--Micha{\l}ek \cite{LM}, Proposition 2.3, with the
one-dimensional restriction. The proof above records the precise
arbitrary-order image form and its norm certificate.
For complex border rank, the bound must hold on the whole complex
image $U_\C$, not just its real points. For example
$\Span_\R\{I_2,\left(\begin{smallmatrix}0&-1\\1&0\end{smallmatrix}\right)\}$
has no nonzero real matrix of rank one, but its complexification does.
Replacing complex minima by real minima would give a false assertion.

\begin{finding}{F-14}{explicit image-norm bounds for the target tensors}
For $M_n$ the first-flattening image has $\alpha(U)=1/n$, and hence
\begin{equation}\label{eq:strong-MM}
 \BR(M_n)\geq n^2+n-1.
\end{equation}
For the $k\mid b$ flattenings, $b=n-k\geq2$, the image constants are
\[
\begin{array}{c|c|c}
T& q=\dim U&\alpha(U)\text{ or a certified upper bound}\\ \hline
D_n&\binom nk&1/b!\\
P_n&\binom nk&b!\binom nb/n^b\\
\mathcal D_n&\binom nk^2&\binom nb/n^b\\
\mathcal P_n&\binom nk^2&\alpha(U)\leq b!\binom nb/n^b .
\end{array}
\]
The first three displayed constants are exact. In particular, putting
$h=\lfloor n/2\rfloor$ and $b=\lceil n/2\rceil$ for $n\geq3$ gives
\begin{align}
 \BR(D_n)&\geq\binom nh+b-1,\label{eq:strong-D}\\
 \BR(P_n)&\geq\binom nh+1,\label{eq:strong-P}\\
 \BR(\mathcal D_n)&\geq
  \binom nh^2+\ceil{\left(\frac{n^b}{\binom nb}\right)^{1/(b-1)}}-1,
        \label{eq:strong-polyD}\\
 \BR(\mathcal P_n)&\geq\binom nh^2+1.\label{eq:strong-polyP}
\end{align}
The polynomial bounds also hold for symmetric border rank.
For example \eqref{eq:strong-polyD} gives
$\BR(\mathcal D_4)\geq38$. All these bounds hold over $\C$ and
for real Euclidean border rank.
\end{finding}
\begin{proof}
Up to row/column permutations and transposition, the image for $M_n$
consists of matrices $I_n\otimes X$. Their operator norm is $\op X$
and their Hilbert norm is $\sqrt n\,\HS X$.
Thus the maximum squared ratio is $1/n$, attained for rank-one $X$.
F-13 with $q=n^2,b=2$ gives \eqref{eq:strong-MM}.

The image for $D_n$ is the alternating subspace $\Lambda^b\F^n$
embedded in the tensor power. Its orthogonal antisymmetrizer satisfies
\[
 \|P_U(x_1\otimes\cdots\otimes x_b)\|^2
     =\frac{\det(\langle x_i,x_j\rangle)_{ij}}{b!}\leq1/b!.
\]
Orthogonal $x_i$ give equality. A stronger uniform matrix-norm
estimate in this image gives \eqref{eq:strong-D}. For a unit
alternating $u$, apply the same unitary in all $b$ factors to
rotate a first-factor test vector to $e_1$.
In an orthonormal wedge basis, the squared Hilbert norm of that
contraction is $1/b$ times the sum of squared wedge coefficients
whose index set contains $1$. It is at most $\HS u^2/b$.
Thus the first matrix flattening of every $u$ has spectral stable
rank at least $b$, and therefore $\BR(u)\geq b$.
Use \eqref{eq:image-sub}.

If $U\subseteq\Sym^b H$, Banach's symmetric spectral theorem B-01
and interchange of compact maxima give
$\alpha(U)=\max_{\|x\|=1}\|P_Ux^{\otimes b}\|^2$.
For $P_n$, an orthonormal basis of $U$ is indexed by $b$-element
subsets, each vector being the normalized sum of its $b!$ orders.
Consequently
$\|P_Ux^{\otimes b}\|^2=b!\,e_b(|x_1|^2,\ldots,|x_n|^2)$,
where $e_b$ is the elementary symmetric polynomial.
For nonnegative entries of sum one,
$e_b\leq\binom nb/n^b$. One direct proof averages two entries:
the only varying term with their sum fixed is their product times
the nonnegative $e_{b-2}$ of the other entries. Repeatedly averaging
a largest and a smallest entry converges to the constant vector:
the variance decreases by half their squared difference.
This process never decreases $e_b$.
Equality at constant moduli proves the claimed exact $\alpha$.

For $\mathcal D_n$, the normalized symmetric tensors of all
$b$ by $b$ minors form an orthonormal basis of $U$, by their
disjoint monomial supports and unit Hilbert norms. For a unit
matrix $X$,
\[
 \|P_UX^{\otimes b}\|^2
 =\sum_{|I|=|J|=b}|\det X_{I,J}|^2
 =e_b(s_1(X)^2,\ldots,s_n(X)^2)
 \leq\binom nb/n^b .
\]
The middle equality follows by applying an SVD to the exterior
power $\Lambda^bX$. Equality holds at $I/\sqrt n$.
For $\mathcal P_n$, $U$ is contained in the symmetric subspace
supported on ordered matrix-unit lists with distinct row indices.
The squared norm of the projection of $X^{\otimes b}$ to that
larger subspace is $b!\,e_b(r_1,\ldots,r_n)$, with
$r_i=\sum_j|X_{ij}|^2$ and $\sum_i r_i=1$.
Projection to $U$ is no larger, proving the stated upper bound.

For $b\geq2$, $b!\binom nb/n^b=\prod_{j=0}^{b-1}(1-j/n)<1$.
F-13 therefore gives an improvement of at least one in both
permanent cases, and gives \eqref{eq:strong-polyD}.
All calculations allowed complex vectors and attainments were real.
\end{proof}

These are norm certificates combined with known border substitution.
They are not asserted to be new rank records.
They do improve a plain flattening and show why an assertion that
all norm-assisted methods are bounded by flattening rank would be
false. The orthogonal projector depends nonlinearly on the image;
the argument also tests incidence with a secant variety. F-02 and
the linear-map barrier below do not apply to this combined argument.
No polynomial-time procedure for arbitrary projector injective norms
or a constant-factor Grothendieck relaxation for them is asserted.

\begin{finding}{F-23}{exact reduction of the residual images}
Let $T_n$ be any of $D_n,P_n,\mathcal D_n,\mathcal P_n$,
and let $2\leq b\leq n-1$. In the $(n-b)\mid b$
flattening image $U_{n,b}$,
\begin{equation}\label{eq:exact-residual}
 \min_{0\ne u\in U_{n,b}}\BR(u)=\BR(T_b).
\end{equation}
Consequently, if $s=1$ for the row families and $s=2$
for the polynomial families,
\begin{equation}\label{eq:recursive-image}
 \BR(T_n)\geq\binom nb^{s}+\BR(T_b)-1.
\end{equation}
These are CP border ranks over $\C$, or real Euclidean border
ranks with the same field on both sides.
In particular, for $n\geq3$,
\[
 \BR(\mathcal D_n),\BR(\mathcal P_n)\geq\binom n2^2+3.
\]
The last inequality holds also at $n=2$ directly. It gives $39$
at $n=4$. Every CP lower bound here also bounds symmetric
border rank when the target is symmetric.
\end{finding}
\begin{proof}
For a row tensor, an image basis is indexed by $b$-element
coordinate subsets $I$; its vectors are the alternating,
respectively symmetric squarefree, tensors on those subsets.
Write a nonzero image vector in that basis and select $I$
with nonzero coefficient. Apply the coordinate projection
$\F^n\to\F^I$ in every remaining factor. Every other
$b$-element subset has a coordinate outside $I$ and vanishes.
The selected vector becomes a nonzero scalar multiple of $D_b$
or $P_b$, with the prescribed normalization.

For a polynomial tensor, the image basis is indexed by pairs
$(I,J)$ of $b$-element row and column subsets. Its vectors
are the normalized coefficient tensors of the corresponding
$b$ by $b$ minors or permanents. Choose a nonzero coefficient
and apply in every factor the coordinate projection
$\F^{n\times n}\to\F^{I\times J}$. A different row or column
subset necessarily uses a matrix entry outside this block
in every one of its monomials, and is annihilated. The
selected term becomes a nonzero multiple of $\mathcal D_b$
or $\mathcal P_b$.

A tensor product of linear maps takes pure tensors to pure
tensors or zero; it therefore cannot increase rank or border
rank. This proves the lower inequality in
\eqref{eq:exact-residual}. Conversely, a basis vector is an
embedded copy of $T_b$. Its rank and border rank equal those
of $T_b$: inclusion gives one inequality and coordinate
projection back gives the other. This proves equality.
F-13 and the image dimensions of F-03 now prove
\eqref{eq:recursive-image}. Finally, the normalized quadratic
coefficient matrices of $\det_2$ and $\operatorname{per}_2$
each consist, after permutation, of two nonzero two-by-two
off-diagonal blocks. Both have matrix rank four. Take $b=2$.
\end{proof}
This is an algebraic refinement of the norm-assisted image
method, with no priority claim. A certified CP lower bound
$l_b$ for the smaller tensor can be substituted on the right
of \eqref{eq:recursive-image}; in particular, the norm bounds
of F-14 can be used recursively. A lower bound for the smaller
tensor's \emph{symmetric} border rank cannot be substituted
unless it is independently known to bound its CP border rank.
For example the known CP equality
$\BR(\mathcal D_3)=17$ from \cite{CHL}, Theorem 1.2,
gives $\BR(\mathcal D_5)\geq116$ and
$\BR(\mathcal D_6)\geq416$.
Conner--Huang--Landsberg \cite{CHuangL}, Theorem 1.1 and
the following identification, give $\BR(\mathcal P_3)=16$;
thus $\BR(\mathcal P_5)\geq115$ and
$\BR(\mathcal P_6)\geq415$.
These complex CP lower bounds also bound real Euclidean
border rank. The minimum-residual gate is thus
an exact smaller-instance problem, not merely an unspecified
projector optimization.

\section{Comparison with stronger known bounds and barriers}
\begin{background}{B-04}{selected comparison theorems}
The following are literature results, not new findings or computations
reperformed here. They also prevent interpreting our norm bounds as
current rank records.
\end{background}

For matrix multiplication, Landsberg--Ottaviani \cite{LO},
Theorem 1.1 and Corollary 1.2, give
$\BR(M_n)\geq2n^2-n$ using restricted Koszul flattenings.
Landsberg--Micha{\l}ek \cite{LM}, Theorem 1.1, give
$2n^2-\ceil{\log_2n}-1$, combining border substitution with
Koszul flattenings. These are comparison formulas, not claims of
best bounds in every small dimension: Conner--Harper--Landsberg
\cite{CHL}, Theorem 1.1, give $\BR(M_3)\geq17$.
F-14 is weaker than the first formula by $(n-1)^2$.
By F-02 any bound that actually comes from one specified linear
Koszul matrix can be rewritten as a nuclear/operator certificate
after matrix preconditioning. This does not eliminate the work of
computing that matrix's rank.

For the row tensors, Han--Ju--Kim \cite{HJK},
Proposition 3.2 and Theorem 4.2, use recursive Koszul maps and prove
\[
 \BR(D_n)\geq\frac{n^{n-1}}{(n-1)!}=\frac{n^n}{n!}
 \quad(n\geq3,\ \text{characteristic zero}).
\]
Their Theorem 5.6 gives $\BR(P_4)=\TR(P_4)=8$ and the lower bounds
$15,29,55$ for $\BR(P_5),\BR(P_6),\BR(P_7)$.
The later \cite{HJK4}, Theorem 1.1, proves $\BR(D_4)=12$
over subfields of $\C$, by border apolarity.
Han--Song \cite{HS}, Theorem A, prove $\TR(P_5)=16$ in
characteristic zero; Proposition 4.1 gives
$\BR(P_n)\geq\binom n{\lfloor n/2\rfloor}+1$ over $\C$ and
its subfields, the same numerical permanent bound as F-14.
The paper explicitly leaves $\BR(P_5)\in\{15,16\}$ unsettled.
Its rank theorem must not be used as a border-rank equality.

The symmetric row permanent $P_n$ represents a scalar multiple of
the monomial $x_1\cdots x_n$. Han \cite{Han}, Theorem A,
proves in a 2026 preprint that its \emph{symmetric} border rank is
$2^{n-1}$ in characteristic zero. This is neither a theorem about
the permanent in $n^2$ variables nor an equality for ordinary
CP border rank.

For the symmetric determinant polynomial tensor, Farnsworth
\cite{Farn}, Theorem 1.3, gives for $n\geq5$
\[
 \BSR(\mathcal D_n)\geq(1+\delta_n)
                   \binom n{\lfloor n/2\rfloor}^{\!2},
\quad
\delta_n=
\begin{cases}
\dfrac{8(n^3+6n^2-8)}{(n-1)(n+2)(n+4)^2(n^2-2)},&n\text{ even},\\[5pt]
\dfrac{16(n^2+8n+9)}{(n+3)(n+5)^2(n^2-2)},&n\text{ odd}.
\end{cases}
\]
Theorem 1.6 gives $\BSR(\mathcal D_4)\geq38$.
Theorem 1.8 and Corollary 1.11 give the historical interval
$14\leq\BSR(\mathcal P_3)\leq16$.
Those are symmetric-border statements; lower bounds on $\BSR$
alone do not lower-bound CP border rank.
The later \cite{CHL}, Theorem 1.2, gives the stronger and different
statement $\BR(\mathcal D_3)=17$ over $\C$.
It is about the tensor in $(\C^9)^{\otimes3}$, not $D_3$.
Conner--Huang--Landsberg \cite{CHuangL}, Theorem 1.1,
settles the permanent case with $\BR(\mathcal P_3)=16$.
Together with the symmetric rank upper bound $16$ from F-04,
this also gives $\BSR(\mathcal P_3)=\SR(\mathcal P_3)=16$.
Our norm argument for $\mathcal D_4$ gives the value $38$ for
ordinary CP border rank as well; F-23 improves it to $39$
for both degree-four polynomial tensors. No priority or
current-record assertion is made.

\begin{background}{B-05}{barriers with a stated scope}
For a nonzero linear map $L:(\C^N)^{\otimes3}\to\operatorname{Mat}_{u,v}$,
put $\rho(L)=\max_{\text{pure }x}\rank L(x)$ and
$\operatorname{Pot}(L)=\max_T\rank L(T)/\rho(L)$.
The classical general rank-method barrier is
$2^dN^{\lfloor d/2\rfloor}$ for order $d$, from
Efremenko--Garg--Oliveira--Wigderson \cite{EGOW},
Theorem 3 and Corollary 4; it gives $8N$ for three factors.
A 2026 preprint of Lovitz \cite{Lov}, Theorem 1.2 and Corollary 1.3,
states the sharper bounds
\[
 2N-2\leq\sup_L\operatorname{Pot}(L)
       \leq2N-2+\frac2{N+1},
\]
with equality at the upper bound for odd $N$.
Accordingly the associated integer border-rank certificates cannot
exceed $2N-1$.
\end{background}
This recent preprint is cited background, not a proof independently
reproduced here. F-02 inherits this ceiling, including arbitrary
invertible output preconditioning. For $M_n$ substitute $N=n^2$.
F-13 is outside this single-linear-map argument because it uses
subspace/secant incidence. Nonlinear equations, border substitution,
and arbitrary GCT obstructions are not ruled out.
For growing tensor order the three-factor ceiling must not be
substituted with $N=n$; one must use the actual order and dimensions,
or the exact matrix sizes and pure-atom costs.

There is also a representation-size issue. A matrix SDP is polynomial
in its explicit matrix size, not necessarily in the parameter $n$
of a succinct tensor formula. Balanced catalecticants and many
Koszul matrices have exponentially many rows or columns in $n$.
Existence of a finite norm certificate is not by itself an efficient
algorithm in that succinct parameter.

\section{The constant, approximation, and what has not changed}
\begin{background}{B-06}{Unique Games threshold}
The exact real Grothendieck constant is
$K_G^\R=\sup_{A\ne0}S(A)/Q(A)$.
Raghavendra--Steurer \cite{RS}, Theorem 1.1, prove that,
assuming UGC, approximation to the real bipartite quadratic/sign
problem with loss factor strictly below $K_G^\R$ is NP-hard.
Thus such a polynomial-time algorithm is excluded under UGC
and $\mathsf P\ne\mathsf{NP}$. Their algorithm obtains arbitrarily close reciprocal
guarantees $1/K_G^\R-\eta$, and Theorem 1.3 makes the constant
computable to any prescribed accuracy in time depending on that
accuracy. For the algorithm's correctly directed bound see
\S3.3, equations (3.5)--(3.7).
\end{background}
The author-hosted PDF prints $K_G-\eta$ in the displayed
Theorem 1.2 under a loss-factor convention, inconsistent with
Theorem 1.1. Its proof gives the reciprocal guarantee just stated,
equivalently loss factor $K_G+\epsilon$.
We use the proof's direction and record this source discrepancy.
Nothing here resolves UGC or identifies the exact constant.

The recent primary sources located during this run report
\[
 1.773\leq K_G^\R<1.77987164,\qquad
 1.4<K_G^\C<1.404898554746.
\]
For the real lower bound see Heilman--Jones--Malavolta
\cite{HJMreal}, Theorem 1.1; for the upper bound see Peng
\cite{Peng}, Theorem 1.1 and \S6, (41)--(43).
The complex interval is \cite{HJMcomplex}, Theorem 1.1, (1.2).
These are date-specific, computer-assisted preprint reports.
Their large numerical certificates have not been rerun here.
We do not label the table an exhaustive global record: Peng
mentions a stronger reported Ren value not independently
source-verified in this run. The exact constants remain unknown;
a numerical interval is not an exact value.
None of F-01--F-14 depends on these new numerical endpoints.

\begin{background}{B-07}{hardness of generic tensor norms}
Friedland--Lim \cite{FL}, Theorem 8.1, give the stated
NP-hardness for real tensor spectral and nuclear norms of order
at least three; their complex result is Theorem 8.10 for order
at least four. The supplied paper must not be cited as proving
its complex three-tensor case. The input model is generic
explicit tensors; it does not prohibit the special closed
formulas in Table \ref{tab:norms}.
\end{background}

\begin{background}{B-08}{cut norm and XOR interpretation}
Alon--Naor \cite{AN}, Lemma 2.1, Theorem 4.3 and \S4.4,
connect the sign problem to cut-norm approximation and its
rounding. Tsirelson's correlation theorem identifies bipartite
real vector correlations with the quantum correlations relevant
to XOR games: Pisier \cite{Pis}, Theorem 12.12 and \S19.
Thus $Q$ and $S$ have classical and quantum bias interpretations.
They are not tensor CP-rank measures.
\end{background}
F-09 certifies the direction of the sign-norm approximation.
F-11 gives a matrix-rank test only with dimension-sensitive
constants and an independent sign upper certificate.
F-05 rules out a fixed-$K_G$ repair of the false full-tensor
rank inequalities. F-13--F-14 add a different, nonlinear
image-norm route, without determining $K_G$ or changing
the UGC threshold.

\part{The findings in Grothendieck's frameworks}
\section{Integral tensor spaces and secant equations}
The translation in this part is mathematical, not a claim that
Grothendieck stated the modern complexity conclusions. The archival
locators in \S\ref{sec:archive} identify the underlying constructions;
the deductions are proved here. Analytic inequalities retain their
chosen archimedean metrics. There is no metric-free version obtained
merely by changing terminology.

\begin{finding}{F-15}{an integral secant model}
Let $V_i=\Z^{n_i}$, $V=\bigotimes_iV_i$, and
$\mathbb A(V)=\Spec\Sym(V^*)$. The universal sum morphism
\[
 \phi_r:\left(\prod_i\mathbb A(V_i)\right)^r\longrightarrow\mathbb A(V),
 \qquad(v_{ij})\longmapsto\sum_{j=1}^r\bigotimes_i v_{ij}
\]
has scheme-theoretic image defined by $\ker\phi_r^*$.
Its coordinate ring is torsion-free, hence flat over $\Z$.
Its complex fiber is the usual affine Segre secant cone.
Every equation in this kernel vanishes on sums of $r$ pure
tensors over every commutative ring.
\end{finding}
\begin{proof}
The coordinates of $\phi_r$ are integer polynomials
$\sum_j\prod_i x_{ij,k_i}$. Its image coordinate ring is the
image of $\Sym(V^*)$ in the polynomial ring of the source.
This is a subring of an integral domain, hence a domain and a
torsion-free $\Z$-module. A torsion-free module over $\Z$ is flat.
The flat base change $\Z\to\C$ preserves the kernel of the
coordinate homomorphism, so the complex image scheme is the
closure of the complex sums. Substitution in an arbitrary ring
preserves each integer polynomial identity.
\end{proof}
Kernel formation need not commute with nonflat reduction modulo
every prime. Thus the statement does not identify every special
fiber with the independently formed secant scheme. The foundation
for the kernel and closure construction is EGA I \cite{EGA1},
\S9.5, Propositions 9.5.1 and 9.5.4, Corollary 9.5.2,
Definition 9.5.3.

\begin{finding}{F-16}{bundle form of the flattening certificates}
Let $A,B,C$ be finite locally free sheaves on a scheme $S$.
The map $K_p(t)$ of F-12 is a canonical morphism
$B^*\otimes\Lambda^p A\to\Lambda^{p+1}A\otimes C$,
compatible with arbitrary base change.
For the universal sum of $r$ pure tensors all
$(r\binom{a-1}p+1)$-minors vanish over $\Z$.
\end{finding}
\begin{proof}
The map is the composition of contraction, tensoring with
$\Lambda^p A$, and exterior multiplication. All are canonical
for locally free sheaves and commute with base change.
In a frame its entries are linear in $t$ with coefficients
$0,1,-1$. F-12 proves the pure-tensor rank calculation over
a field. Apply its rank bound to the universal sum over $\Q$;
the specified minors are zero polynomials over $\Q$.
They have integer coefficients, so are zero over $\Z$ and
after every base change.
\end{proof}
For any such matrix of bundle maps, rank at most $s$ is expressed
by $\Lambda^{s+1}\Phi=0$, equivalently its $(s+1)$-minors.
With target rank $f$, these minors generate
$\operatorname{Fitt}_{f-s-1}(\coker\Phi)$ when that index is meaningful.
At an archimedean point, F-02 supplies an analytic certificate
that some such minor does not vanish. The minor, not a
Hermitian norm over $\Z$, is the algebraic information.
Exterior powers and their $K_0$ classes are treated in
SGA 6 \cite{SGA6}, Expos\'e V, Examples 2.2(b) and Lemma 2.2.1.

\section{Hilbert schemes, limiting spans, and image incidence}
\begin{finding}{F-17}{the FGA parameter space retaining a limiting span}
Let $X\subset\CP(V)$ be integral, nondegenerate, and projective
over an algebraically closed field. For $1\leq r\leq\dim V$,
let $U_r\subset\Hilb^r(X)$ parametrize reduced length-$r$
subschemes of independent points. Close the graph of its span map
$U_r\to\Gr(r,V)$ in $\Hilb^r(X)\times\Gr(r,V)$, obtaining $\Gamma$.
If $\mathcal S$ is the pullback of the tautological rank-$r$
subbundle, the image of
$\CP(\mathcal S)\to\CP(V)$ is exactly $\sigma_r(X)$.
\end{finding}
\begin{proof}
The Hilbert scheme and Grassmannian are projective, hence so
are $\Gamma$ and the projective bundle; its image is closed.
Over $U_r$ the image consists of the spans of independent
distinct $r$-tuples. Such tuples form a nonempty dense open
of $X^r$: nondegeneracy allows their successive choice, and
$X^r$ is irreducible. The sum of the tautological lines over
this dense open parametrizes a dense part of the usual secant
construction. The corresponding open of $\CP(\mathcal S)$
is dense because $\Gamma$ is the graph closure.
Continuity puts the image of the whole projective bundle in
the closure of the open image; properness gives the reverse
containment. This is $\sigma_r(X)$.
\end{proof}
Here projectivization parametrizes lines. This construction
retains both the length-$r$ scheme and its limiting $r$-plane.
It does not identify that plane with the actual span of the
limiting scheme if the evaluation rank drops. At the boundary
$Z\subset\CP(L)$, but equality of spans is not automatic.
This distinction is why border rank is not casually replaced
by a statement about the actual span of an arbitrary limiting
Hilbert point. The representability input is FGA \cite{FGA},
Bourbaki expos\'e 221, Theorems 3.1--3.2 and the Hilbert
specialization after Proposition 3.8.
Projectivity implies properness by EGA II \cite{EGA2},
Theorem 5.5.3(i); universal closedness and base change are
Definition 5.4.1 and Proposition 5.4.2.

\begin{finding}{F-18}{the scheme dictionary for F-13}
Let $Y\subset\CP(W)$ be the Segre variety for the remaining
tensor factors. For $s\geq1$ the locus
\[
 \mathcal I_{q,s}
 =\{U\in\Gr(q,W):\CP(U)\cap\sigma_s(Y)\ne\varnothing\}
\]
is closed. If a complex tensor has a flattening image of dimension $q$
and border rank at most $r$, that image belongs to
$\mathcal I_{q,r-q+1}$.
\end{finding}
\begin{proof}
In $\Gr(q,W)\times\sigma_s(Y)$ take the incidence of pairs
$(U,[w])$ with $w\in U$. It is the zero locus of the
tautological line's map into the universal quotient bundle.
Projection to the Grassmannian is proper, since the second
factor is projective; its image is closed and is precisely
$\mathcal I_{q,s}$.
The approximating $q$-planes constructed in F-13 each meet
$\sigma_{r-q+1}(Y)$ by the atom-span dimension argument.
Closedness puts their limit in the same incidence locus.
\end{proof}
Thus a uniform image-norm estimate in F-14 certifies avoidance
of a projective secant-incidence locus. The projector is metric
data on the complex tautological bundle. Avoidance is the
algebraic-geometric conclusion. The projectors do not themselves
form regular algebraic functions over $\Z$.
The universal incidence, exterior-power, and proper-image
constructions make sense in relative form over $\Z$; as in F-15,
care is needed in identifying special-fiber secants.
For a Waring-only problem the residual atom variety must be
Veronese; it cannot silently be replaced by Segre.
Our lower bounds use the larger CP atom set and consequently
also bound symmetric border rank.

\section{Group schemes, \texorpdfstring{$K_0$}{K0}, and what a class forgets}
The group scheme $\prod_i\GL(V_i)$ acts on the tensor bundle.
Explicitly $\GL_n=\Spec\Z[x_{ij},\det(x_{ij})^{-1}]$.
For an integral tensor $t$, the kernel of the orbit-map
homomorphism $\Sym(V^*)\to\Z[\prod_i\GL(V_i)]$ defines an
integral orbit-closure model; flat base change to $\C$ gives
the usual complex orbit closure, as in F-15.
SGA 3 \cite{SGA3}, Expos\'e I, \S4.5--\S4.6, supplies
the group and module schemes.
In characteristic zero the polynomial problem instead uses
$\GL(W)$ on $\Sym^n(W^*)$ and the Veronese variety; integral
symmetrization requires the convention in F-22 below.

\begin{finding}{F-19}{metrics are extra data}
No nonzero norm on a tensor space is invariant under the full
product-$\GL$ action. Even the ratios in \eqref{eq:ratios}
generally vary along one orbit.
\end{finding}
\begin{proof}
Scalar multiplication by $\lambda>1$ in one factor would imply
$N(\lambda T)=N(T)$ under invariance, whereas homogeneity gives
$N(\lambda T)=\lambda N(T)$.
For the ratios take
$T_t=e_1^{\otimes d}+t e_2^{\otimes d}$, $0<t\leq1$, $d\geq2$.
All lie in one product-$\GL$ orbit. Their norms are
$\HS{T_t}^2=1+t^2$, $\eps(T_t)=1$, $\pi(T_t)=1+t$.
For the spectral claim use Cauchy--Schwarz on the two coordinate
products in the first two factors and bound the other coordinates
by one; $e_1^{\otimes d}$ attains one.
The displayed decomposition proves the nuclear upper bound;
pairing with $e_1^{\otimes d}+e_2^{\otimes d}$, of spectral norm
one, proves its lower bound. Thus $b(T_t)=(1+t)^2/(1+t^2)$ varies.
\end{proof}
The exact dictionary for F-01--F-14 is therefore a locally free
sheaf with a chosen Euclidean/Hermitian metric on an archimedean
fiber, not just its underlying scheme. Norm duality, convexity,
and PSD statements remain real/complex analytic statements.
Optimizing a certificate over an orbit is an additional operation,
not a consequence of writing $\GL$ as a group scheme.

For finite locally free sheaves,
$\lambda^i[E]=[\Lambda^iE]$ in $K_0(S)$, and exterior-power
filtrations make $\lambda_t([E]+[F])=\lambda_t[E]\lambda_t[F]$.
For representations use $K_0$ of the exact tensor category of
finite locally free representations. Over $\C$ for a reductive
group this recovers the usual representation ring.
The Koszul source and target have classes
$[B^*]\lambda^p[A]$ and $[C]\lambda^{p+1}[A]$.
Schur functors give the corresponding associated bundles and
characters. Over $\Z$ one uses integral Schur functors;
Young idempotents involving factorial or hook-length division
are only available after the corresponding localization.
These constructions explain dimensions and representation types,
not the rank of a particular section-induced map.

\begin{finding}{F-20}{a precise $K_0$ and $\gamma$ limitation}
The source and target classes of a morphism, even with all their
$\lambda$- and $\gamma$-operations, do not determine its rank.
The virtual class of its two-term complex also does not do so.
\end{finding}
\begin{proof}
Over $\Spec\C$, compare the zero and identity maps
$\C^N\to\C^N$. Both have the same source/target class pair,
and all ring, $\lambda$, or $\gamma$ expressions in those
classes therefore agree. Their ranks are $0,N$.
The virtual two-term classes are both $[\C^N]-[\C^N]=0$.
In the family $z\,\id$ over $\mathbb A^1$, these classes remain
constant although rank jumps at $z=0$.
For an extension example in $\operatorname{Rep}_\C(\mathbb G_a)$,
the trivial two-dimensional representation and
$z\mapsto\left(\begin{smallmatrix}1&z\\0&1\end{smallmatrix}\right)$
both have class $2[\mathbf1]$, but only the first has trivial action.
\end{proof}
This does not say that $K_0$ of an image, or tensor-dependent
equivariant constructions, is useless. Those may contain the
rank information one first needs to determine.
The exact $\lambda$ and $\gamma$ foundations are SGA 6,
Expos\'e V, Definition 2.1, Examples 2.2(b), Lemma 2.2.1,
the representation-ring example, Proposition 3.2(ii), and
\S3.10, (3.10.1). The last defines the filtration by products
of $\gamma^{i_j}(x_j)$ with augmentation-zero $x_j$ and
$\sum i_j$ at least the filtration degree.

\section{Weights, purity, and motives: the exact nonimplication}
\begin{finding}{F-21}{constant stalks and absolute point motives do not give rank}
Over $\mathbb F_q$, with $\ell\ne\operatorname{char}\mathbb F_q$,
the zero tensor point and a nonzero pure tensor point have the
same constant-$\Q_\ell$ stalk with Frobenius weight zero, and
the same absolute point motive, but different tensor ranks.
\end{finding}
\begin{proof}
Each rational point as an abstract scheme is $\Spec\mathbb F_q$.
The constant sheaf's stalk and degree-zero cohomology are
$\Q_\ell$, with Frobenius eigenvalue one. Both are pure of
weight zero. Both abstract points have the unit motive.
Their ranks in the ambient tensor representation are zero and one.
\end{proof}
This rules out only the proposed inference from these specified
data. Sheaves constructed from secant geometry or relative motives
with the embedding retained may contain further information;
none is constructed by the norm proofs here.
The \'etale site is SGA 4 \cite{SGA4}, Expos\'e VII, \S1.2.
Purity is used in the sense of Deligne \cite{WeilII},
Definition 1.2.2. No theorem from SGA 5, or the weights theorem
of Weil II, is asserted to imply a tensor-rank lower bound.
There is consequently no fabricated SGA 5 or manuscript locator
for such a nonexistent implication.

\section{Integral determinant and permanent, without dividing by \texorpdfstring{$n!$}{n!}}
\begin{finding}{F-22}{an integral catalecticant dictionary}
The determinant and permanent are integral polynomials in
$\Sym^n(E^*)$, $E=\Z^{n^2}$. Their normalized Hilbert tensors
$\mathcal D_n,\mathcal P_n$ are naturally defined over
$\Z[1/n!]$, not as integral sections with the same coefficients.
The bidegree-$(k,n-k)$ coproduct of $F(x+y)$ defines an integral map
\[
 C_k(F):\Gamma^k E\longrightarrow\Sym^{n-k}(E^*).
\]
For determinant and permanent this map has rank $\binom nk^2$
over every field. On a polynomial power $\ell^n$ its rank is
at most one, possibly zero in small characteristic.
\end{finding}
\begin{proof}
A squarefree degree-$n$ monomial has $n!$ ordered tensor
positions, so the coefficient of each normalized symmetric
position is its polynomial coefficient divided by $n!$.
The integral coproduct avoids that division:
$\Gamma^kE$ is the dual of $\Sym^k(E^*)$, and its columns
are the Hasse derivatives
$\partial^{[\alpha]}x^\beta=\binom\beta\alpha x^{\beta-\alpha}$.
For determinant/permanent, a nonzero column differentiates a
matching of $k$ distinct rows and columns. It equals, up to
determinant sign, the complementary minor or permanent.
Every pair $(I,J)$ of $k$-element row and column sets occurs.
The $\binom nk^2$ complementary polynomials have disjoint
monomial supports and coefficients $\pm1$ or $1$, so remain
nonzero and independent over every field. This gives their
rank and a free integral basis for the image.
Finally the relevant coproduct of a power is
$\binom nk\ell(x)^k\ell(y)^{n-k}$, an outer product,
so has rank at most one after every base change.
\end{proof}
Hence its $(r+1)$-minors vanish on sums of $r$ polynomial
powers and their closures, as a universal integral assertion.
In characteristic zero it is the catalecticant in F-03 up to
the conventional nonzero basis scalings.
In characteristic $p\leq n$, the polynomial power map
$\ell\mapsto\ell^n$ and the complete-linear-system Veronese
model must not be identified through characteristic-zero
symmetrization: multinomial coefficients can vanish.
Divided powers provide the appropriate integral distinction.
The no-rank-drop statement above is specific to this Hasse
matrix; arbitrary flattening ranks can drop modulo primes.
None of these integral assertions equips the fibers with a
Hilbert norm.

\section{Orbit-closure obstructions and the representation ring}
\begin{background}{B-09}{the no-occurrence theorem, with its range}
Let $Z_{n,m}$ be the padded-permanent orbit closure and
$\Omega_n$ the determinant orbit closure in the degree-$n$
polynomial representation. B\"urgisser--Ikenmeyer--Panova
\cite{BIP}, Theorem 1.4 in arXiv v3, prove that for
$n\geq m^{25}$ every irreducible occurrence in
$\C[Z_{n,m}]$ also occurs in $\C[\Omega_n]$.
Their Corollary 2.6 gives the analogous comparison with a
power-sum orbit closure.
\end{background}
An actual equivariant closed inclusion $Z\hookrightarrow\Omega$
induces a graded equivariant surjection of coordinate rings.
Over $\C$ semisimplicity therefore forces multiplicity inequalities
in every degree. Occurrence-support inclusion is weaker:
it retains only whether a coefficient of an irreducible class in
the representation ring is zero. $K_0$ retains multiplicities;
passing to support discards more information.
Thus B-09 is not orbit-closure inclusion, not a prohibition of
multiplicity obstructions, and not a consequence of a tensor-norm
barrier. The supplied Landsberg survey \cite{Land}, \S2.2--\S2.3,
provides the Segre/Veronese and equation dictionary; its \S2.9
concerns a different nonnormality issue.

\section{Exact archival and textual crosswalk}\label{sec:archive}
The \emph{R\'esum\'e} is catalogue item \textbf{CSG17} in the
CSG published-works archive. That catalogue dates it 1953;
Grothendieck Circle's bibliography also catalogues it as 1956.
We do not infer an unverified manuscript/publication chronology
from those differing dates. The facsimile used is the 1996 USP
reprint, DOI \texttt{10.11606/resimeusp.v2i4.74836}, with its
original internal pagination and 80 PDF leaves.
The table gives \emph{one-based} scan leaves; subtract one for
a zero-based PDF index. A leaf printed $8'$ explains why the
offset is not constant throughout the facsimile.

\begingroup\small
\begin{longtable}{p{2.1cm}p{7.7cm}p{1.5cm}p{1.5cm}}
\toprule
Source key & Exact textual locator & Printed page & Scan leaf\\\midrule
\endhead
G53-N & CSG17, \S1 no.1, Definition 1 and Theorem 1:
extremal reasonable norms and projective infimum formula & 8 & 8\\
G53-D & CSG17, \S1 no.1, Theorem 2 and duality corollary & $8'$ & 9\\
G53-G & CSG17, \S4 no.2, Theorems 1 and 2, fundamental theorem & 59--60 & 60--61\\
G53-I & CSG17, \S4 no.2, Corollary 1 after Theorem 3,
inequality (8) & 61 & 62\\
G53-P & CSG17, \S4 no.3, opening of the proof & 62 & 63\\
EGA-image & EGA I, 9.5.1--9.5.2, and 9.5.3--9.5.4;
NUMDAM \texttt{PMIHES\_1960\_\_4\_\_5\_0} & 176--177 & 175--176\\
EGA-proper & EGA II, Definition 5.4.1, Proposition 5.4.2;
Theorem 5.5.3(i); NUMDAM \texttt{PMIHES\_1961\_\_8\_\_5\_0}
& 100; 104 & 97; 101\\
SGA3-G & SGA 3 I, \S4.5--\S4.6, general linear and module schemes & 29 & 42\\
SGA6-L & SGA 6 V, Definition 2.1; Examples 2.2(b),
Lemma 2.2.1 & 314--316 & 321--323\\
SGA6-R & SGA 6 V, representation-ring example following the
exterior-power discussion & 317 & 324\\
SGA6-G & SGA 6 V, Proposition 3.2(ii);
\S3.10, (3.10.1), $\gamma$-filtration & 325; 331 & 332; 338\\
FGA-Q & Bourbaki expos\'e 221, Theorems 3.1--3.2;
NUMDAM \texttt{SB\_1960-1961\_\_6\_\_249\_0}, catalogue CSG44
& 259--260 & 12--13\\
FGA-H & Same expos\'e, Hilbert specialization after Proposition 3.8;
Remarks 3.9 & 265--266 & 18--19\\
WII & Deligne, Weil II, Definition 1.2.2, in the IAS facsimile
listed in \cite{WeilII} & 153 & 17\\
\bottomrule
\end{longtable}
\endgroup
The original R\'esum\'e, EGA, and SGA pages listed here were
visually inspected by the geometry agent; the check ledger
distinguishes source-statement/location checks from proof audits.
The modern \texttt{AG/AG-22.pdf} transcription was only a
cross-check, not substituted for the scan.
\textbf{CSG17 is a published-work catalogue identifier, not a
Montpellier unpublished-manuscript shelfmark. No such manuscript
cote was established or is claimed.} The verified published
locator and facsimile are supplied explicitly rather than
inventing an additional archival field.

\subsection{Every finding and its exact dictionary}
\begingroup\small
\begin{longtable}{p{1.8cm}p{9.0cm}p{2.1cm}}
\toprule
Finding IDs & Reformulation and proof transport & Text anchors\\\midrule
\endhead
F-01 & On a metrized archimedean fiber, mode-image subspaces
are supports of the tensor section; their dimensions are
determinantal ranks. Orthogonal expansion proves the same
inequalities; F-15--F-16 carry the rank conclusion integrally.
& G53-N/D; EGA-image\\
F-02, F-03 & Bundle maps have nonvanishing exterior minors.
Hilbert duality certifies their ranks; F-22 supplies the
integral polynomial catalecticant. The singular values require
the stated metrics.
& G53-D; SGA6-L; EGA-image\\
F-04 & Finite signed averages are identities of symmetric
tensor sections over characteristic zero. Their coefficient
proofs survive algebraically after removing denominators;
nuclear lengths are archimedean data.
& G53-N/D; SGA3-G; F-22\\
F-05 & Integral tensor identities for multiplication and
decompositions give the rank upper bounds. Tensor products
and reindexing are morphisms of locally free modules.
The analytic ratios on their real/complex fibers disprove
the proposed rank conversion.
& G53-N/D; SGA3-G\\
F-06 & $T_t=W+t(\cdots)+t^2e_2^{\otimes3}$ is a polynomial
family over $\Z[t]$; over $t\ne0$ it is a sum of two powers.
Its special fiber lies in a secant closure. The slice pencil
proves its rank is three. Norms remain fiberwise.
& EGA-image; FGA-H; G53-D\\
F-07, F-08 & Compact normalized atoms and norm duality live
on the metrized fiber. Unrestricted scheme-theoretic
degenerations do not impose bounded coefficients. The
identity-matrix counterexamples are sections of free bundles.
& G53-N/D; FGA-H\\
F-09--F-11 & PSD Gram matrices and real sign cubes belong
to the analytic real fiber. Grothendieck's inequality
compares the specified Banach tensor norms, not an
\'etale or $K_0$ norm. Any certified matrix-rank conclusion
is a determinantal condition on the corresponding bundle map.
& G53-G/I/P; EGA-image\\
F-12, F-16 & Contraction followed by exterior multiplication
is a natural morphism of locally free sheaves. Its
pure-section rank calculation proves universal integral
minor equations.
& SGA6-L; SGA3-G\\
F-13, F-14, F-18, F-23 & The image is a point of a Grassmannian;
border rank forces incidence with a residual Segre secant.
Projector injective norms certify avoidance on the complex
fiber. The proper incidence proof transports the limiting
argument exactly. Coordinate quotients of free sheaves isolate
a smaller determinant or permanent section and prove the
exact residual-rank reduction.
& FGA-Q/H; EGA-proper; G53-N/D\\
F-15, F-17 & Scheme-theoretic images and projective
Hilbert/Grassmann incidence express secant membership.
Retaining the limiting plane prevents an unjustified
span-of-limit identification.
& EGA-image/proper; FGA-Q/H\\
F-19 & The same group scheme preserves rank but does
not preserve a chosen norm. Scalar multiplication and
the diagonal family prove the failure directly.
& SGA3-G; G53-N\\
F-20 & Classes of sources, targets, and two-term complexes
coincide for maps of different rank. All $\lambda$ and
$\gamma$ expressions in those classes coincide as well.
The morphism itself must be retained.
& SGA6-L/R/G\\
F-21 & Two abstract rational points have identical
constant \'etale stalks, weight zero, and unit motives.
The embedding carries the differing tensor rank.
No tensor-rank theorem is imported from purity.
& SGA 4 VII 1.2; WII\\
F-22 & The coproduct of an integral symmetric polynomial
uses the divided-power dual domain. Primitive complementary
minor supports give a characteristic-independent rank;
normalized symmetric Hilbert tensors require $1/n!$.
& SGA6-L; SGA3-G; EGA-image\\
\bottomrule
\end{longtable}
\endgroup
For relative image Grassmannians over a possibly nonreduced base,
one works on the determinantal rank-$q$ stratum: the
$(q+1)$-minors vanish scheme-theoretically and the $q$-minors
generate the unit ideal. An invertible $q$-minor then puts the
map locally in block form with an identity block and zero
complementary block, so its image is a subbundle.
Constant ranks on geometric fibers alone would not suffice.

\part{Unchecked findings, remaining gates, and implications}
\section{Status at session close}
There are \textbf{no unchecked new mathematical assertions
presented as findings}. F-01--F-23 have full proofs and
recorded whole-proof model checks. The computational checks
in the package are finite exact identities and normalization
checks, not formal verification of the entire paper.
Cited background is not relabelled as an independently
reproved result. In particular, the large computations behind
the recent numerical constants and cited border-apolarity or
recursive-Koszul results were not rerun.

The following stronger claims are not established:
\begin{enumerate}
\item An improved asymptotic matrix-multiplication exponent
lower bound, a superpolynomial determinantal-complexity lower
bound, or a separation of $\mathsf P$ and $\mathsf{NP}$.
Our proved rank inequalities and counterexamples do not
imply any of them.
\item A universal impossibility theorem for norm-based,
semidefinite, nonlinear, or geometric methods.
F-02 and B-05 concern specified linear-map certificates;
F-13 is an explicit reason not to overstate that barrier.
\item A polynomial-time way to compute the residual
projector norm in F-13 for arbitrary input, or a reduction
of that norm to the classical Grothendieck SDP with a
dimension-independent guarantee.
\item An exact value of either Grothendieck constant, or
an independently rerun certification of the 2026 endpoints.
\item A tensor-dependent \'etale sheaf or motive with a
proved weight-to-rank implication. F-21 only refutes a
specific automatic translation using constant-point data.
\item An unpublished-manuscript archival cote for the
R\'esum\'e. The published-work catalogue ID, exact original
facsimile leaves, EGA/SGA/FGA locators, and the distinction
between those kinds of reference are verified.
\end{enumerate}
The first five are limits of the proved scope, not missing
lemmas in any stated finding. The last is a bibliographic
field left explicitly unverified. If an unpublished-manuscript
cote is required in addition to the published archive
locator, that archival subrequest remains incomplete.

\section{Implications for the complexity question}
The useful positive result is more precise than a choice
between norms and geometry. Nuclear duality certifies
matrix ranks of flattenings; an injective norm of the
\emph{entire image projector}, combined with border
substitution, can improve a plain flattening. The improvements
in F-14 follow with explicit constants and complex-field
proofs; F-23 reduces the residual problem for all four
determinant/permanent families exactly to a smaller member.
The latter mechanism depends on nonlinear image and
incidence data and has no generic computational efficiency
claim. It does not change the Grothendieck constant.

The determinant/permanent norm table also shows why exact
large norms alone do not separate arithmetic complexity:
both determinant and permanent tensors have large projective
norms, and the determinant polynomial itself has exponential
catalecticant ranks despite efficient determinant evaluation.
Waring/CP rank measures restricted decompositions, not
unrestricted arithmetic circuits or determinantal complexity.
Orbit-closure noncontainment requires equations or
representation-theoretic data for the relevant polynomial
orbits. The present inequalities do not produce such an
obstruction. B-09 remains a restriction on occurrence
obstructions in its stated range, not on every such approach.

If the next objective is a stronger GCT bound, the concrete
gate left by this work is a stronger uniform certificate for
the residual image's avoidance of a secant variety, or a
different tensor-dependent construction that escapes the
identified linear-map ceiling. Merely replacing a norm by a
constant-factor SDP approximation, or a bundle map by its
$K_0$ source and target classes, does not pass that gate.
\texttt{LEDGER.md} records the attempted routes, corrections,
and this endpoint; \texttt{QUESTIONS.tsv} records each
question's proved, refuted, or remaining status.

\begin{thebibliography}{99}
\bibitem{Gro}
A. Grothendieck, \emph{R\'esum\'e de la th\'eorie m\'etrique
des produits tensoriels topologiques}, catalogue CSG17,
1953 (also catalogued as 1956), Boletim Soc. Mat. S\~ao Paulo
8, 1--79; facsimile reprint 1996.
\url{https://revistas.usp.br/resenhasimeusp/en/article/download/74836/78406/100978}.
DOI \texttt{10.11606/resimeusp.v2i4.74836}.
Catalogue: \url{https://csg.igrothendieck.org/published-works/}.
Exact locators and scan leaves: \S\ref{sec:archive}.

\bibitem{FL}
S. Friedland and L.-H. Lim, \emph{Nuclear Norm of Higher-Order Tensors},
arXiv:1410.6072v3, 18 May 2016 (the supplied source).
\url{https://arxiv.org/pdf/1410.6072}.
Used locators: \S2 (5)--(8); Proposition 3.1; \S5 (23),
Theorem 5.1, Corollaries 5.2--5.3; Lemma 6.2 (34)--(37);
Theorems 8.1, 8.10; Lemma 9.2 (57), display before (60),
Theorem 9.4. The source corrections are F-08.

\bibitem{Der}
H. Derksen, \emph{On the Nuclear Norm and the Singular Value
Decomposition of Tensors}, arXiv:1308.3860v2, 21 April 2014.
\url{https://arxiv.org/pdf/1308.3860}.
Corollaries 1.2, 1.10; Theorems 1.14--1.15; \S1.5 (4)--(5);
\S1.7 (8); \S8, p.~23, unnumbered five-term determinant
decomposition and final display before Remark 8.1.

\bibitem{CLL}
E. Carlen, E. H. Lieb, M. Loss,
\emph{An Inequality of Hadamard Type for Permanents},
arXiv:math/0508096v1, Theorem 1.1, equation (1.1).
\url{https://arxiv.org/pdf/math/0508096}.

\bibitem{Pis}
G. Pisier, \emph{Grothendieck's Theorem, past and present},
arXiv:1101.4195v3, 4 June 2011; supplied
\texttt{grothendieck.bams-june3.tex}.
\url{https://arxiv.org/pdf/1101.4195}.
Theorems 2.4, 3.1, 12.12; equations (2.4)--(2.5),
(3.13)--(3.14), (4.4); \S\S19, 22.

\bibitem{AN}
N. Alon and A. Naor, \emph{Approximating the Cut-Norm via
Grothendieck's Inequality}, SIAM J. Comput. 35 (2006), 787--803.
Lemma 2.1, Theorem 4.3, \S4.4.
\url{https://web.math.princeton.edu/~nalon/PDFS/cutnorm3.pdf}.

\bibitem{RS}
P. Raghavendra and D. Steurer, \emph{Towards Computing the
Grothendieck Constant}, SODA 2009, 525--534.
Theorems 1.1, 1.3; algorithm proof \S3.3, (3.5)--(3.7).
\url{https://www.dsteurer.org/paper/grothendieck.pdf}.
The displayed Theorem 1.2 direction is discussed in the text.

\bibitem{HJMreal}
S. Heilman, C. Jones, G. Malavolta,
\emph{Primal/Dual Method for the Grothendieck Constant},
arXiv:2610.10477v2, 14 November 2026, Theorem 1.1;
certification components Theorems 4.1--4.2.
\url{https://arxiv.org/pdf/2610.10477}.

\bibitem{Peng}
B. Peng, \emph{Certified Fourier rounding for the real
Grothendieck inequality}, arXiv:2610.24417, October 2026.
Theorem 1.1; \S6 (41)--(43).
\url{https://arxiv.org/abs/2610.24417}.
Author-posted full text inspected at
\url{https://hub.valency.io/works/W_uhkyapjd}.

\bibitem{HJMcomplex}
S. Heilman, C. Jones, G. Malavolta,
\emph{Sharper Bounds for the Complex Grothendieck Constant},
arXiv:2609.25462, September 2026, Theorem 1.1 (1.2);
Theorem 7.1.
\url{https://www.stevenheilman.org/~heilman/papers/grothcx.pdf}.

\bibitem{LO}
J. M. Landsberg and G. Ottaviani,
\emph{New Lower Bounds for the Border Rank of Matrix Multiplication},
Theory of Computing 11(11) (2015), 285--298.
Theorem 1.1, Corollary 1.2, p.~286; Theorem 2.1.
\url{https://theoryofcomputing.org/articles/v011a011/v011a011.pdf}.

\bibitem{LM}
J. M. Landsberg and M. Micha{\l}ek,
\emph{A $2n^2-\log_2(n)-1$ lower bound for the border rank
of matrix multiplication}, arXiv:1608.07486v1;
IMRN 2018(15), 4722--4733.
Theorem 1.1; Proposition 2.3, p.~3; \S2, (2)--(4).
\url{https://arxiv.org/pdf/1608.07486}.

\bibitem{EGOW}
K. Efremenko, A. Garg, R. Oliveira, A. Wigderson,
\emph{Barriers for Rank Methods in Arithmetic Complexity},
ITCS 2018, LIPIcs 94, Article 1.
Theorem 3, p.~1:10; Corollary 4, p.~1:11.
\url{https://www.math.ias.edu/~avi/PUBLICATIONS/EfremenkoGaOlWi2018.pdf}.

\bibitem{Lov}
B. Lovitz, \emph{Explicit tensors beyond the linear flattening barrier},
arXiv:2610.08504v1, October 2026, Theorem 1.2 and
Corollary 1.3, pp.~2--3.
\url{https://arxiv.org/abs/2610.08504}.
Full manuscript inspected at \url{https://borowiecki.dev/pdf/2610.08504}.
Preprint statement used as background.

\bibitem{HJK}
J. I. Han, J.-H. Ju, Y. Kim,
\emph{Recursive Koszul flattenings of determinant and permanent tensors},
arXiv:2503.12032v1, 15 March 2025.
Proposition 3.2; Theorem 4.2, p.~9; Theorem 5.6, pp.~14--15.
\url{https://arxiv.org/pdf/2503.12032}.

\bibitem{HJK4}
J. I. Han, J.-H. Ju, Y. Kim,
\emph{The border rank of the $4\times4$ determinant tensor is twelve},
arXiv:2510.11051v2, 18 August 2026, Theorem 1.1, p.~2.
\url{https://arxiv.org/pdf/2510.11051}.

\bibitem{Han}
J. I. Han, \emph{The border Waring rank of $x_1\cdots x_n$
is $2^{n-1}$}, arXiv:2609.09141v1, 8 September 2026,
Theorem A, p.~2.
\url{https://arxiv.org/pdf/2609.09141}.

\bibitem{HS}
J. I. Han and J. Song, \emph{The rank of the $5\times5$
permanent tensor is sixteen}, arXiv:2608.09708v2,
27 August 2026, Theorem A, p.~3; Proposition 4.1, pp.~10--11.
\url{https://arxiv.org/pdf/2608.09708}.

\bibitem{Farn}
C. Farnsworth, \emph{Koszul--Young Flattenings and Symmetric
Border Rank of the Determinant}, arXiv:1505.05079v2,
11 November 2015. Theorems 1.3, 1.6, 1.8; Corollary 1.11.
\url{https://arxiv.org/pdf/1505.05079}.

\bibitem{CHuangL}
A. Conner, H. Huang, J. M. Landsberg,
\emph{Bad and good news for Strassen's laser method:
Border rank of $\operatorname{perm}_3$ and strict submultiplicativity},
arXiv:2009.11391v1, 23 September 2020.
Theorem 1.1 and the following identification, p.~3; proof in \S5.
\url{https://arxiv.org/pdf/2009.11391}.

\bibitem{CHL}
A. Conner, A. Harper, J. M. Landsberg,
\emph{New lower bounds for matrix multiplication and $\det_3$},
Forum of Mathematics, Pi 11 (2023), e17,
DOI \texttt{10.1017/fmp.2023.14}.
Theorems 1.1--1.2, \S1.2, p.~4.
\url{https://arxiv.org/abs/1911.07981}.

\bibitem{BIP}
P. B\"urgisser, C. Ikenmeyer, G. Panova,
\emph{No occurrence obstructions in geometric complexity theory},
arXiv:1604.06431v3, 17 September 2018 (supplied source).
Theorem 1.4, p.~3; Corollary 2.6, p.~7.
\url{https://arxiv.org/pdf/1604.06431}.
Older version numbering is not used.

\bibitem{Land}
J. M. Landsberg, \emph{Geometric complexity theory: an introduction
for geometers}, supplied \texttt{gctsurvey12-10.tex}.
\S2.2, Relevant algebraic varieties; \S2.3, First equations;
\S2.9, Bad news. These are definitions and context, not the
source of the later BIP theorem.

\bibitem{EGA1}
A. Grothendieck and J. Dieudonn\'e,
\emph{\'El\'ements de g\'eom\'etrie alg\'ebrique I},
Publ. Math. IH\'ES 4 (1960), \S9.5, pp.~176--177.
NUMDAM \texttt{PMIHES\_1960\_\_4\_\_5\_0}.
\url{https://www.numdam.org/article/PMIHES_1960__4__5_0.pdf}.

\bibitem{EGA2}
A. Grothendieck, with J. Dieudonn\'e,
\emph{\'El\'ements de g\'eom\'etrie alg\'ebrique II},
Publ. Math. IH\'ES 8 (1961).
Definition 5.4.1 and Proposition 5.4.2, p.~100;
Theorem 5.5.3(i), p.~104.
\url{https://www.numdam.org/article/PMIHES_1961__8__5_0.pdf}.

\bibitem{SGA3}
M. Demazure and A. Grothendieck et al.,
\emph{Sch\'emas en groupes}, SGA 3, Expos\'e I,
\S4.5--\S4.6, p.~29.
\url{https://library.slmath.org/nonmsri/sga/sga/pdf/sga3-1.pdf}.

\bibitem{SGA6}
P. Berthelot, A. Grothendieck, L. Illusie et al.,
\emph{Th\'eorie des intersections et th\'eor\`eme de Riemann--Roch},
SGA 6, Expos\'e V, Definition 2.1, Examples 2.2(b),
Lemma 2.2.1, representation-ring example, Proposition 3.2(ii),
\S3.10, (3.10.1).
\url{https://library.slmath.org/nonmsri/sga/sga/pdf/sga6.pdf}.

\bibitem{FGA}
A. Grothendieck, \emph{Techniques de construction et th\'eor\`emes
d'existence en g\'eom\'etrie alg\'ebrique IV: les sch\'emas de Hilbert},
S\'eminaire Bourbaki, expos\'e 221 (May 1961), 249--276.
Theorems 3.1--3.2; after Proposition 3.8; Remarks 3.9.
NUMDAM \texttt{SB\_1960-1961\_\_6\_\_249\_0}, catalogue CSG44.
\url{https://www.numdam.org/article/SB_1960-1961__6__249_0.pdf}.

\bibitem{SGA4}
M. Artin, A. Grothendieck, J.-L. Verdier et al.,
\emph{Th\'eorie des topos et cohomologie \'etale des sch\'emas},
SGA 4, Expos\'e VII, \S1.2 (the \'etale site).
No tensor-rank implication is attributed to this construction.

\bibitem{WeilII}
P. Deligne, \emph{La conjecture de Weil II},
Publ. Math. IH\'ES 52 (1980), 137--252.
Definition 1.2.2, p.~153.
\url{https://publications.ias.edu/sites/default/files/Number40.pdf}.
\end{thebibliography}
\end{document}





